\section{A full-dimensional seed component for the actual cubic tuple}
\label{sec:cubic-tuple-seed-component}

The native support ring and the invariant deformation ring have
different roles.  Finiteness and torsionfreeness over the native
weight algebra do not pass automatically to a quotient of the
invariant ring.  We instead construct a surjection from that ring
onto one particular native component.  This is the dimension
argument required by the literal \path{SeedIn} predicate.

\paragraph{Actual data and retained inputs.}
Retain the actual arithmetic quotient, local groups and
normalization element of AC-G94--AC-G99, with \(p=269\).
Thus \(\mathcal O\) has finite residue field \(k\),
\(G\) is the actual joint-kernel quotient on the CM field,
\(g^\star=cg^{-1}c^{-1}\), and
\[
 \boldsymbol\theta=(\chi^j\gamma,1,\chi^j\gamma^{-1}),\qquad
 w=(1,0,-1),\qquad
 a_0=\theta_0(g_0)-\theta_2(g_0)\ne0.
 \tag{G107.1}
\]
In particular \(\theta_i(g^\star)=\theta_{2-i}(g)\).
The local separating elements and the arithmetic
\(\Phi\)-condition are those already verified in AC-G94.
Let \(R\), \(B\subset A\subset R\), and \(r^{\rm univ}\)
be the actual tuple ring and algebras of AC-G94--AC-G95.

Write \(\mathcal G_0\) for the native unramified-outside-\(S\)
CM Galois group and
\[
 \pi:\mathcal G_0\twoheadrightarrow G,\qquad
 \kappa=\epsilon^{-1}
 \tag{G107.2}
\]
for AC-G97.  The native residual characters are
\(\bar t_i=(\theta_i\pi)\bar\kappa\).
The kernel of \(\pi\) is the closed normal,
conjugation-stable subgroup generated by the native
\(\operatorname{Ip}(u)\) and \(I_0\).
On all of these inertia groups \(\kappa=1\).
The full coefficient-prime local groups and their maps are
the same on both sides of (G107.2).

Let \(\mathscr B\) be the actual native universal ordinary
determinant ring for this residual datum.  Its existence,
the native \(\mathrm{ModU}\) input, and the actual arithmetic
seed construction remain precisely the retained source inputs
of AC-G89 and their host extensions in AC-G91/G93.
Use AC-G97's map
\[
                  \Phi:\mathscr B\longrightarrow B
 \tag{G107.3}
\]
and let \(J\subset\mathscr B\) be its native support ideal.
We use the following conclusions already supplied there:
\[
 \begin{gathered}
 \mathscr B/J\ne0,\qquad
 \mathscr B/J\text{ is finite and torsionfree over }
 \Lambda\simeq\mathcal O[[X_1,\ldots,X_{3d}]],\\
 \mathrm{Unr0}(J),\\
 \phi:\mathscr B\longrightarrow K',\quad J\subset\ker\phi,
 \quad K'\text{ algebraically closed of characteristic zero},\\
 D_{\phi}=\det(\nu),\qquad
 \nu:\mathcal G_0\longrightarrow\operatorname{GL}_3(K')
       \text{ absolutely irreducible}.
 \end{gathered}
 \tag{G107.4}
\]
Here \(\nu\) is the actual native normalization of the
automorphic seed; its determinant and ordered local characters
are those of the native pair at \(\phi\).
No original single-character \(\mathrm{TwLink}\) or
\(\mathrm{NRep.IsBar}\) is substituted for the actual tuple.

\begin{lemma}[AC-G107-integral: integral invariants of a tuple field point]
\label{AC-G107-integral}
Let \(C\) be a complete Noetherian local
\(\mathcal O\)-algebra which is a domain and has residue field
\(k\), and let \(K\) be an algebraically closed field containing
\(C\).  Suppose that an absolutely irreducible representation
\(\rho:G\to\operatorname{GL}_3(K)\) has a continuous trace
\(t:G\to C\), with
\[
 \operatorname{tr}\rho(g)=t(g),\qquad
 \bar t(g)=\sum_i\theta_i(g),\qquad
 t(g^\star)=t(g).
 \tag{G107.5}
\]
Suppose also that continuous characters
\(\zeta_{v,i}:G_v\to C^\times\), with the actual ordered
residual characters, satisfy
\[
 \prod_{i=0}^2
 \bigl(\rho(\operatorname{loc}_v h_i)
                      -\zeta_{v,i}(h_i)I\bigr)=0
       \quad(h_i\in G_v).
 \tag{G107.6}
\]
Then, after conjugating \(\rho\) over \(K\), there are maps
\[
 \xi:A\longrightarrow K,\qquad \beta:B\longrightarrow C
 \tag{G107.7}
\]
of \(\mathcal O\)-algebras such that the representation and
ordered characters at \(\xi\) are exactly
\((\rho,\zeta)\), and
\[
                   \xi|_B=\iota_C\beta .
 \tag{G107.8}
\]
The conjugated representation satisfies the actual polarity
and normalization equations of AC-G94.  This assertion
does not require a \(C\)-free stable lattice for \(\rho\).
\end{lemma}

\begin{proof}
\emph{1. Polarity and the integral normalization scalar.}
The representation \(g\mapsto\rho(g^\star)^t\) has the
same trace as \(\rho\).
Absolute irreducibility and the trace pairing give an
intertwining invertible matrix \(P\).
For completeness, the equality of traces identifies the
kernels of the two maps from \(K[G]\) onto \(M_3(K)\):
if one image of \(x\) is zero, pair the other image with
every group matrix and use the nondegenerate matrix trace
pairing.  Thus the induced map of matrix algebras is an
automorphism.  It is inner: the images of the coordinate
idempotents give three complementary lines, and rescaling
their chosen bases makes the images of all matrix units
the usual matrix units.  This proves the asserted
intertwiner without a semisimplicity assertion about any
integral lattice.

The usual transpose argument, or Schur's lemma applied to
the two intertwining forms, gives \(P^t=\lambda P\),
where \(\lambda^2=1\).
An invertible alternating matrix cannot have odd rank in
characteristic different from two.  Hence \(P\) is symmetric,
and we may write the relation as
\[
       \rho(g^\star)=P^{-1}\rho(g)^tP.
 \tag{G107.9}
\]
Put \(Q=\rho(g_0)-\rho(g_0^\star)\).
It is skew-adjoint for \(P\).
In rank three its characteristic polynomial is
\[
 \operatorname{charpol}_Q(X)=X(X^2-b),\qquad
 b=\tfrac12\operatorname{tr}(Q^2)
   =t(g_0^2)-t(g_0g_0^\star)\in C.
 \tag{G107.10}
\]
Indeed skew-adjointness gives \(\operatorname{tr}Q=0\)
and \(\det Q=0\); the remaining coefficient follows by
the degree-two Newton identity.  These calculations
use only that two is a unit.
Reduction of the trace expression gives
\[
             \bar b=
        \bigl(\theta_0(g_0)-\theta_2(g_0)\bigr)^2=a_0^2.
 \tag{G107.11}
\]
This uses the reversal identity for the tuple at \(g_0\);
it does not assume that the endpoint product is one.
Simple-root Hensel lifting in \(C\) supplies the unique
\(a\in C^\times\) with \(a^2=b\) and \(\bar a=a_0\).

The three eigenvalues \(a,0,-a\) of \(Q\) are distinct.
Their eigenspaces pair under \(P\) only when their
eigenvalues sum to zero.  The two outer lines therefore
pair nondegenerately with each other, and the middle
line has nonzero self-pairing.
Rescale these bases, taking a square root for the middle
line in \(K\), to put the form into the antidiagonal
matrix \(J_3\).
In this basis
\[
 \rho(g^\star)=J_3\rho(g)^tJ_3,\qquad
 Q=\operatorname{diag}(a,0,-a).
 \tag{G107.12}
\]
The traces and the ordered-product equations are
unchanged by this constant conjugation.

\emph{2. The finite trace order and continuity.}
Let
\[
                  \mathcal E=\sum_{g\in G}C\rho(g)
                      \ \subset M_3(K).
 \tag{G107.13}
\]
This is a unital \(C\)-algebra.
Choose nine group matrices \(\rho(g_s)\) forming a
\(K\)-basis of \(M_3(K)\), which is possible by absolute
irreducibility over the algebraically closed field.
The Gram matrix
\[
             \bigl(t(g_sg_t)\bigr)_{s,t}
 \tag{G107.14}
\]
has nonzero determinant \(d_0\in C\).
Trace coordinates show that
\(\mathcal E\subset d_0^{-1}\sum_s C\rho(g_s)\).
The containing module is finite over the Noetherian ring
\(C\), so \(\mathcal E\) is finite as well.
Its trace lies in \(C\).

The map
\[
 \mathcal E\longrightarrow C^9,\qquad
 X\longmapsto
       \bigl(\operatorname{tr}(X\rho(g_s))\bigr)_s
 \tag{G107.15}
\]
is an injective map of finite \(C\)-modules.
Artin--Rees identifies its induced maximal-ideal topology
with the module topology on \(\mathcal E\).
Since \(g\mapsto t(gg_s)\) is continuous for every \(s\),
\(g\mapsto\rho(g)\) is continuous with values in this
finite module.  This supplies the continuity which is
not part of an abstract field-valued determinant realization.

Write
\[
       M_{ij}=\sum_{g\in G}C\rho(g)_{ij}\subset K.
 \tag{G107.16}
\]
Each is a finite \(C\)-module, being an image of
\(\mathcal E\), and each entry map is continuous for
its module topology.

\emph{3. Diagonal entries and balanced products.}
The spectral projectors of \(Q\) are
\[
 E_0=\frac{Q(Q+aI)}{2a^2},\qquad
 E_2=\frac{Q(Q-aI)}{2a^2},\qquad E_1=I-E_0-E_2.
 \tag{G107.17}
\]
They belong to \(\mathcal E\), and in the normalized
basis are precisely the coordinate projectors.
Thus
\(\rho(g)_{ii}=\operatorname{tr}(E_i\rho(g))\in C\).
To check the residues and the cycle products without
assuming an integral matrix lattice, lift the displayed
polynomials in \(Q\) to \(C[G]\).
The \(C\)-linear trace on that group algebra reduces
to the trace of
\(\bar\rho=\operatorname{diag}(\theta_0,\theta_1,\theta_2)\),
by (G107.5).
The reductions of the lifted \(E_i\) are the coordinate
projectors, because the residual difference at \(g_0\)
is \(\operatorname{diag}(a_0,0,-a_0)\).
Consequently
\[
 \overline{\rho(g)_{ii}}=\theta_i(g),\qquad
 \rho(g_1)_{i_0i_1}\cdots
 \rho(g_s)_{i_{s-1}i_0}\in\mathfrak m_C
 \tag{G107.18}
\]
whenever the displayed cycle includes an off-diagonal
factor.  The second product is the trace of the
corresponding product of projectors and group matrices;
its residual trace is zero.

Give an entry weight \(w_i-w_j\).
Every weight-zero monomial which contains an
off-diagonal entry belongs to \(\mathfrak m_C\).
Here is the full rank-three verification.
Remove diagonal factors, which already lie in \(C\).
A balanced multiset of weights in \(\{1,2,-1,-2\}\)
decomposes into the patterns
\[
 (1,-1),\quad(2,-2),\quad(1,1,-2),\quad(2,-1,-1).
 \tag{G107.19}
\]
Remove opposite equal weights first; if both signs
remain, the lighter weight occurs twice against the
heavier one.  This proves the decomposition by induction.
Polarity in (G107.12) identifies the two positions
of weight \(1\), and likewise the two of weight \(-1\),
by replacing a group argument with its star.
Therefore the four patterns can respectively be
written as cycles with index lists
\[
        (0,1,0),\quad(0,2,0),\quad
        (0,1,2,0),\quad(0,2,1,0).
 \tag{G107.20}
\]
The entries are scalars in the field \(K\), so
reordering their factors is legitimate.
Equation (G107.18) proves the claim for each factor
and hence for the original balanced monomial.

\emph{4. An actual complete test ring and its invariant image.}
With an indeterminate \(z\), form the graded subring
\[
 D=C[\,z^{w_i-w_j}\rho(g)_{ij}:g\in G,\ i,j\,]
                           \ \subset K[z,z^{-1}].
 \tag{G107.21}
\]
Finite generation of the modules (G107.16) shows that
\(D\) is a finite-type \(C\)-algebra.
It is a domain, and the preceding balanced-product
calculation shows that \(D_0=C\).
For \(f\in D\), let \(f_0\) be its Laurent coefficient
of degree zero.  Then
\[
          D\longrightarrow k,\qquad
             f\longmapsto f_0\bmod\mathfrak m_C
 \tag{G107.22}
\]
is a surjective ring homomorphism.
Indeed a product of two components of nonzero,
opposite degrees is a sum of balanced monomials with
an off-diagonal entry, so its constant coefficient
is in \(\mathfrak m_C\).  This proves multiplicativity;
additivity and the assertion on constants are literal.
Let \(\mathfrak n\) be its kernel, and set
\[
                         \widehat D=\widehat{D_{\mathfrak n}}.
 \tag{G107.23}
\]
This is a complete Noetherian local
\(\mathcal O\)-algebra with residue field \(k\).
The maps \(D\to D_{\mathfrak n}\to\widehat D\) are
injective: the first uses that \(D\) is a domain,
and the second is Krull intersection for the
Noetherian local ring \(D_{\mathfrak n}\).
Thus \(C\), \(D\), and their displayed relations may
all be regarded as actual subrings of \(\widehat D\).

Consider the matrices
\[
 \rho^z(g)=
 \operatorname{diag}(z,1,z^{-1})\rho(g)
 \operatorname{diag}(z^{-1},1,z).
 \tag{G107.24}
\]
Their entries and those of their inverse matrices lie
in \(D\), by applying the same formula to \(g^{-1}\).
Together with the characters \(\zeta_{v,i}\in C^\times\),
they reduce under (G107.22) to the diagonal tuple.
The star, normalization and ordered-product equations
are preserved by this diagonal conjugation.
Continuity into \(\widehat D\) follows from Step 2:
for fixed \(i,j\), the \(C\)-linear map
\(M_{ij}\to\widehat D\), \(x\mapsto z^{w_i-w_j}x\),
sends \(\mathfrak m_C^nM_{ij}\) into
\(\mathfrak m_{\widehat D}^n\widehat D\).
There are only finitely many entry positions.
The local characters are already continuous in \(C\).
These are therefore actual test-ring data for AC-G94,
which supplies a local continuous map
\[
                         R\longrightarrow\widehat D.
 \tag{G107.25}
\]

The image of every zero-weight entry monomial lies in
\(C\), and the same is true of each local character.
The image of \(C\) in \(\widehat D\) is closed.
To see this, \(C/\mathfrak m_C^n\) is finite for every
\(n\), so completeness makes \(C\) compact; its
continuous image in the Hausdorff ring \(\widehat D\)
is closed.  Taking the defining closure of the
weight-zero algebra in AC-G95 consequently puts the
image of all of \(B\) in \(C\).
This defines \(\beta:B\to C\).
Since \(A\) is generated over \(B\) by the actual
matrix entries, its image in (G107.25) lies in \(D\).
Injectivity of \(D\to\widehat D\) makes this an
actual homomorphism \(A\to D\).
Evaluate the Laurent variable at \(z=1\) to obtain
\(\xi:A\to K\).  Equations (G107.7)--(G107.8) and
the representation and character assertions follow.
\end{proof}

\begin{theorem}[AC-G107-seed: transfer of a full native component]
\label{AC-G107-seed}
Under the actual data (G107.1)--(G107.4), set
\(I=\Phi(J)B\).
There are prime ideals \(\mathfrak x\subset A\) and
\(Z\subset B\) with
\[
 \begin{gathered}
 r^{\rm univ}\bmod\mathfrak x\text{ absolutely irreducible},\\
 I\subset Z=\mathfrak x\cap B,\qquad
                         \dim B/Z=3d+1.
 \end{gathered}
 \tag{G107.26}
\]
More precisely, choose a minimal prime
\(Z_{\rm nat}\) over \(J\) contained in \(\ker\phi\).
There is a surjective map
\[
 \beta:B\twoheadrightarrow C=\mathscr B/Z_{\rm nat}
 \quad\text{such that}\quad
           \beta\Phi=q:\mathscr B\twoheadrightarrow C,
 \tag{G107.27}
\]
and \(Z=\ker\beta\).
Thus (G107.26) is the actual changed-tuple
\(\mathrm{SeedIn}(I)\) assertion.
It uses neither torsionfreeness of \(B/I\) nor a claim
that every native component survives that quotient.
\end{theorem}

\begin{proof}
\emph{1. Choose the native component and its generic representation.}
Since \(J\subset\ker\phi\), a minimal prime
\(Z_{\rm nat}\) over \(J\) can be chosen below
\(\ker\phi\).
The quotient \(C=\mathscr B/Z_{\rm nat}\) is a
complete Noetherian local domain with the same
residue field \(k\).
It has characteristic zero, since its map to \(K'\)
extends the given characteristic-zero coefficient map.
The finiteness and torsionfreeness in (G107.4) give
\[
                           \dim C=3d+1.
 \tag{G107.28}
\]
Explicitly, a minimal prime of the finite
\(\Lambda\)-algebra \(\mathscr B/J\) is associated.
No nonzero element of the domain \(\Lambda\) is a
zero-divisor on this algebra, so that prime contracts
to zero.  Its quotient is finite integral over
\(\Lambda\), which preserves dimension.
This is the same exact use of native torsionfreeness
as in AC-G89 and the source \path{dimOf}.

Put \(K=\overline{\operatorname{Frac}C}\).
Base-change the native pair to \(C\) and then to \(K\).
The source theorem
\path{groupDet_exists_semisimple_representation} supplies
a semisimple representation
\(\nu_C:\mathcal G_0\to\operatorname{GL}_3(K)\)
with exactly this determinant law.
This general determinant-realization theorem is
retained as an input, with its algebraically closed
field hypothesis; no residual irreducibility is required.

The representation \(\nu_C\) is absolutely irreducible.
Indeed choose nine group matrices spanning \(M_3(K')\)
at the actual irreducible seed \(\nu\).
Their trace Gram determinant is the specialization
of an element of \(C\) and is nonzero at that seed.
It is therefore nonzero in \(C\), and remains so in
\(K\).  The nine generic matrices span \(M_3(K)\).
This proves generic irreducibility on the chosen
component; irreducibility is not inferred just from
the dimension of that component.

\emph{2. Descend the generic representation through the actual kernel.}
Untwist by \(\kappa\), writing
\[
 \widetilde\rho(g)=\kappa(g)^{-1}\nu_C(g),\qquad
 \widetilde t(g)=\kappa(g)^{-1}\operatorname{tr}\nu_C(g)
                                  \in C.
 \tag{G107.29}
\]
Its group matrices still span \(M_3(K)\).
The trace is continuous for the maximal-ideal topology
of \(C\), by the native pair and continuity of \(\kappa\).
Native polarity, together with
\(\kappa(g)\kappa(cgc^{-1})=\epsilon(g)^{-2}\),
gives
\(\widetilde t(g^\star)=\widetilde t(g)\).

For \(g\in\operatorname{Ip}(u)\), the field
\(\mathrm{Pair.I}\) says
\[
        \operatorname{tr}\bigl(\nu_C(h)(\nu_C(g)-I)\bigr)=0
                           \quad(h\in\mathcal G_0).
 \tag{G107.30}
\]
Full matrix span makes \(\nu_C(g)=I\).
Similarly \(\mathrm{Unr0}(J)\), inherited by \(C\),
implies \(\nu_C(g)=I\) for every \(g\in I_0\).
The character \(\kappa\) is one on these groups,
so \(\widetilde\rho\) kills them too.

Its kernel is closed even before matrix continuity
has been established.  By the trace pairing,
\[
 g\in\ker\widetilde\rho
 \quad\Longleftrightarrow\quad
 \widetilde t(hg)=\widetilde t(h)
                         \text{ for all }h\in\mathcal G_0.
 \tag{G107.31}
\]
The right side is an intersection of closed sets.
It is conjugation-stable under \(c\):
star invariance of the trace and (G107.31) show that
the star of a kernel element is again in the kernel;
taking inverses gives its conjugate by \(c\).
It is already a normal subgroup because it is a
representation kernel.
The exact AC-G97 generation of \(\ker\pi\)
therefore gives
\[
                   \widetilde\rho=\rho\circ\pi
                       \quad\text{for a representation on }G .
 \tag{G107.32}
\]
This uses the actual joint residual-kernel clause.
There is no single-\(\Psi\) substitution.

The trace descends to a continuous \(C\)-valued
function \(t\) on \(G\): a continuous surjection
between these profinite groups is a quotient map.
Its residue is \(\sum_i\theta_i\), and it is
star-invariant.
Divide the native ordered local characters by
\(\kappa\) to get the continuous
\(\zeta_{v,i}:G_v\to C^\times\) with their actual
ordered residual values.
The field \(\mathrm{Pair.Ord}\) says that the trace
of the ordered product, multiplied by every group
matrix, is zero.
Removing the nonzero scalar \(\kappa\)-factors
and using full matrix span gives exactly (G107.6).
The full local groups are unchanged in AC-G97,
so this is the required equation on each actual
local group, not an equation only on its image.

\emph{3. Recover the native component from the invariant ring.}
Apply AC-G107-integral to obtain \(\xi\) and \(\beta\).
Over \(K\), the actual pair associated to the
representation at \(\xi\), after twisting by \(\kappa\),
is the original generic native pair.
The two maps
\[
                  \beta\Phi,\ q:\mathscr B\longrightarrow C
 \tag{G107.33}
\]
therefore carry the universal determinant and
ordered local characters to the same pair over \(K\).
The inclusion \(C\hookrightarrow K\) is injective,
so they carry them to the same pair over \(C\).
For the determinant law this can be checked
coefficientwise on every finite group-algebra support;
it requires no flat-base-change assumption.
For local characters it is ordinary injectivity.

Both maps in (G107.33) are allowed maps to the native
test object \(C\).
In particular every \(\mathcal O\)-algebra map between
the complete local rings with this identified residue
field is local and continuous, by AC-G94's unit argument.
Uniqueness in the native universal property now gives
\(\beta\Phi=q\).
As \(q\) is surjective, \(\beta\) is surjective:
if \(c=q(b)\), then \(c=\beta(\Phi(b))\).
This deduction does not need surjectivity of \(\Phi\).

Set \(\mathfrak x=\ker\xi\) and \(Z=\ker\beta\).
They are prime, since their targets are domains.
Equation (G107.8) and injectivity of \(C\hookrightarrow K\)
give \(Z=\mathfrak x\cap B\).
Equation (G107.27) puts \(\Phi(J)B\) in \(Z\), and
the induced isomorphism \(B/Z\simeq C\), together
with (G107.28), proves the dimension assertion.
The representation at \(\mathfrak x\) becomes the
absolutely irreducible \(\rho\) over \(K\); hence it
is absolutely irreducible over its own fraction
field as well.
In the literal ideal definition it therefore avoids
both triangular ideals and the odd-weight ideal.
Either triangular condition preserves a coordinate
flag, and vanishing of all odd-weight entries
preserves the middle line and endpoint plane.
Each would contradict irreducibility.
This proves exactly the Type I clause of
\(\mathrm{SeedIn}\).

Finally the original native seed is still visible
on this component.  The factorization
\(\phi=\phi_C q\) gives the invariant seed map
\(\phi_C\beta:B\to K'\), which kills \(Z\) and
has composite \(\phi\) with \(\Phi\).
Thus the chosen full-dimensional invariant component
lies under an actual invariant lift of that same
native arithmetic seed.
\end{proof}

\paragraph{Exact source interface.}
The literal definitions are
\path{PropagationTheoremGL2.lean}, \path{SeedIn}, and
\path{SecondThirdPrimeConditionTransferAlt.lean},
\path{Unr0}, \path{DimT},
\path{IdealSpecializationPoint}, and
\path{CurveSpecializationHypothesis}.
Their conclusion requires a Type I prime of \(A\)
and a prime of \(B\) of dimension \(3d+1\) containing
the support ideal and lying below its contraction.
Equation (G107.26) supplies these clauses directly.

The complete source proofs \path{SeedOf.lean},
\path{OfTrace.lean}, \path{OfForm.lean},
\path{PointField.lean}, and \path{AlgVring.lean}
were read.  Their old residual clauses use powers
of one character.  The normalization, balanced-cycle,
topology and test-ring argument above supplies
their changed-tuple counterpart instead of invoking
those clauses unchanged.
The generic containing proof
\path{GroupDetExistsSemisimpleRepresentation.lean},
the quotient proof \path{CoQuot.lean}, and the
native dimension proof \path{DimOf.lean} were also read;
the determinant-realization theorem's deeper imported
inputs are retained.
All source files are at the pinned commit
\path{7cbf7a0c32fca1badf05931879a37c3f8a2b03b4}.
The commutative-algebra inputs are Artin--Rees
(Stacks tag \href{https://stacks.math.columbia.edu/tag/00IN}{00IN}),
Krull intersection
(\href{https://stacks.math.columbia.edu/tag/00IP}{00IP}),
and Noetherian local completion
(\href{https://stacks.math.columbia.edu/tag/0BNH}{0BNH},
Lemmas 10.97.4--10.97.6).
The applied printed proofs were checked; no Lean
execution or new proof of the imported automorphy
statements is claimed.

\paragraph{Scope and geometric dictionary.}
The result closes the actual full-dimensional seed
support clause under the stated native support inputs.
It is stronger than nonemptiness of the characteristic-zero
support: the specified component is recovered as a
quotient of \(B\), with its dimension unchanged.
It does not assert that \(B/I\) is torsionfree over
\(\Lambda\), that all of its components have that
dimension, or that an independently prescribed
characteristic-\(p\) curve lies in the support.
The remaining propagation and numerical hypotheses
must still be applied in their actual proved form.

Geometrically, the trace projectors label the three
residual summands of the local system.
The auxiliary Laurent variable performs diagonal
conjugation; its vertex completion is an actual
formal deformation chart with the required diagonal
specialization.
The induced invariant morphism recovers a whole
native support component.  This uses the existing
formal-deformation and finite-cover dictionary and
asserts no new archival SGA attribution.
